\section{Background}
\subsection{Research Problem}
We\footnote{
{\bf NOTE:} this section contains the complete formal definition of
all approximation approaches. Later, needs to split into two sections:
one for background/walk-around the general idea, and then another section
of algorithms after the preliminary section.}
consider the problem of {\em discharging} a {\em benign} input
string from a large collection of regex malware signatures (i.e.,
the string does not match any of the signatures).

\begin{definition} Let $R = \{r_1, ..., r_n\}$ be a collection
of regex signatures. 
We {\em discharge} a string $s$ for $R$ if 
we show that $\forall i\in [n]: s \not\in L(r_i)$.
\end{definition}

Constructing an efficient zk-prover, even when $|R|=1$ is 
challenging. It is known \cite{AlgBook}[Chapter 5.4] that
a regex $r$ can be translated to NFA size of $3|r|$.
To prove a string $s \not\in L(r)$, one can discharge it
by performing a BFS search over the NFA graph, which results
in complexity $O(|r| |s|)$, and this zk-prover complexity is achieved 
in Zombie \cite{Zombie23}. The concern in practice is that
the concrete prover cost is high. 
The ClamAV logical signatures have an average length of $300$,
% main.ldb size: 11649598/38888
thus to discharge an $1MB$ file against one regex signature needs
over $2^{30}$ R1CS constraints. 

The straw-man approach of discharging regex signatures one by one
will not work, as usually the collection of signatures is large.
For instance, the ClamAV logical regex signatures has
$38k$ patterns. Another approach is to build an automaton
from $\bigwedge_{i=1}^n \bar{r_i}$. This is not practical as
NFA conversion to DFA is needed for regular expression
negation, and DFA intersection causes exponential state explosion.

It may seem completely impractical to produce concretely
efficient zk-proofs for discharging benign files for
malware signature sets used by the industry. 
However, in this paper, we present technical solutions
which allow us to discharge an entire image of binary executables
of a Linux operating system, or the entire Enron email dataset
for all ClamAV signatures and SNORT signatures.
We develop {\em conservative} over-approximations of
regex signatures for discharging a benign input string,
and we design efficient encoding of these approximations
by leveraging folding and IVC schemes in zk-provers.

Both ClamAV and SNORT ruleset support full fledged
PCRE regex patterns. However, not all advanced features of
PCRE regex are used ``equally", which influence the zk-prover design.
Thus, we delve into the technical details of these virus
signatures first, and explore the language features (mainly PCRE)
that we need to support.


\subsection{ClamAV Logical Signatures}
As ClamAV signatures subsume the syntax of SNORT, for introducing
the regex syntax, we focus on
ClamAV only. There are three types of ClamAV signatures: (1) hash 
signatures, where a SHA256 or MD5 is generated directly from a malicious
file; (2) fixed word signatures where Aho-Corasick can be directly
applied to discharge them (and the zk-prover in \cite{Zkreg}
can discharge an entire Linux image efficiently); and (3)
logical regex signatures which empowers the full PCRE regex
syntax. In ClamAV signature set\footnote{ClamAV 0.103.11.27152, retrieved
January 12, 2024}, there are $38839$ signatures. To discharge hash
signatures, one can simply encode the hash operation of an input string
and produce the hash as the secret output witness of the circuit,
and then use a lookup argument (e.g., \cite{Caulk22,CaulkPlus,Baloo,cq})
or zero knowledge non-membership proof to discharge it. The cost is
linear on the input string length and can be independent of the
hash signature set size. In summary, the $38k$ 
logical regex signatures are the ones to discharge in this paper.
We illustrate its syntax using a factious example to integrate
all language features, and for convenience of presentation the 
patterns in the example are much shorter than the real 
signatures.


\begin{example} \label{examp:fake012}
The vast majority of ClamAV signatures are encoded using
4-bit nibbles, because in many cases, binary (executable) contents
are the target. However, ClamAV can be also used to handle HTTP and
email contents. In this case, the text information is represented by
their encoding (e.g., ASCII). For instance an ASCII string
``\texttt{ab}" is encoded as ``\texttt{6162}" for \texttt{0x61} is
the ASCII value for `\texttt{a}'.
The following is a sample ClamAV logical signature.

\begin{small}
\begin{Verbatim}[breaklines=true, breakanywhere=true]
FakeVirus-012;Engine:51-255,Target:0;0|(1=0&2>2);300:616161??3131*31;626262{2-}323232::i;/^([^a]+bc){1,3}/
\end{Verbatim}
\end{small}
\end{example}

A logical signature consists of three parts: (1) the meta information
(e.g., the virus name, the ClamAV engine version needed,
and the \texttt{Target} which specifies the file types such as
executables and emails that should
be checked, and a $0$ indicates all files. \footnote{In this work,
we apply all signatures to all file types as to check file content
type needs expressiveness greater than regex.}. (2) the logical 
relation between sub-signatures. For instance: in Example \ref{examp:fake012},
a file is regarded as containing virus \texttt{FakeVirus-012} if
there is at least one match of sub-signature $0$, or sub-signature $2$
has more than $2$ matches, \footnote{We interpret the semantics as
each match is not nested in another. For instance, for regex
\texttt{a+b}, string \text{aabb} contains $1$ match.}
and sub-signature 1 has no match. Notice that here by specifying
counter constraints \texttt{1=0}, it is essentially specifying the
negation of a regex constraint. (3) sub-signatures. They are
separated by semicolons. There are two types of
sub-signatures: {\em hex} and {\em PCRE} sub-signatures.

A {\em hex} sub-signature is pattern oriented, and is {\em weaker}
then regex, in the way that repetitions (such as Kleene starts)
can be applied to wildcard match only. 
Here a $\verb|?|$ stands for a wild card of one hex nibbles, and
$\verb|{4}|$ stands for $4$ bytes (8 nibbles). 
Note that as a {\em hex} sub-signature
can have wildcard matches, and can express distances between
sub-patterns, it is more expressive than the fixed-word
signatures in ClamAV that directly encoded using Aho-Corasick automata.

A {\em PCRE}
regex allows full PCRE regex syntax such as character classes
such as \texttt{[a-z0-9]}, negated char classes such as \texttt{[\^{}a]},
macros like \texttt{\textbackslash w},
zero-length assertions such as \texttt{\${}},
and also named back-references and look-arounds such as
\texttt{(?P=tagname)}.  
% check code/zkregplus/data/clamav/categories/pcre.dat
Out of the $37955$ logical signatures, there are $559$ ($1.4\%$)
patterns that carry PCRE regex signatures. 
% check data/zkregplus/data/clamav/new_src/README
Note that look-arounds exceed the capability of regex,
and they are approximated. But such cases are rare,
only $6$ out of $559$ PCRE sub-signatures have such features,
and they can be approximated without raising false-positives
of virus on the data set of all Linux binary executables and
Enron emails.
\footnote{
For instance given pattern 
\texttt{<(?P<tag>([a-z]+)).*?>.*?</(?P=tag)>},
we approximate it to: \texttt{<(([a-z]+)).*?>.*?</([a-z]+)>}
in the sense that any concrete string that match the original
regex also matches the over-approximated regex.
There are $6$ patterns (1 positive look-ahead and 5 negative look-ahead)
that use look-around, and we manually convert them to equivalent ones (for 
positive look-ahead) and approximated ones (for negative look-ahead)
for engineering convenience. 
we re-ran ClamAV over data samples
so that no false-positives of virus report are generated.
}

We use a uniform alphabet of hex digits to encode all finite state machines
used in this paper. Thus, the PCRE regex sub-signatures are converted to
this alphabet. One benefit is that the complete DFA resulted from such signatures
needs less space for encoding transitions, because each state has at most
$16$ out-going transitions.  Given the discussion above, we can translate
all sub-signatures of \texttt{FakeVirus-012} into the following over
hex alphabet:

\begin{small}
$r_0$ := \texttt{.\{600\}616161..3131.*31.*}
\end{small}

\begin{small}
$r_1 := \texttt{.*(62|42)(62|42)(62|42).....*323232.*}$
\end{small}

\begin{small}
$r_2 := \texttt{.*((00|...|60|62|...|ff)+6263)\{1,3\}.*}$
\end{small}

Observe $r_0$ for sub-signature $0$: 
it first encodes the position requirement
\texttt{300:}, which requires that the match starts
at exactly the $300$'th byte from the beginning of the
document, into a wildcard of $600$ nibbles. Then each
wildcard match \texttt{?} is translated into a \texttt{.},
and lastly a wildcard string \texttt{.*} is appended at the
end because we need to express the concept the the
pattern is {\em contained} in the target file.

For $r_1$, the ignore case \texttt{::i} requirement is
handled by syntax re-writing. It is apparent that a requirement
of ``sub-signature 1 has no matches" in target string $s$
can be expressed by: $s \in \overline{r_1}$, because the wildcard
strings at the beginning and end of $r_1$.
Similarly, negated char class in $r_2$ is handled by
re-writing. Finally, the logical signature can be
denoted by the following extended regex pattern (in the sense
that intersection and negation are allowed in regex specification):
\[
	r_0 \vee (\overline{r_1} \wedge r_2r_2r_2)
\]

If one wishes to discharge a string $s$ of virus \texttt{FakeVirus-012},
then it is to prove:
\[
	s \in \overline{r_0} \wedge (\overline{r_1} \vee \overline{r_2r_2r_2})
\]
Intuitively, it implies that one has to discharge $s$ 
of $r_1$ first, and then discharge either of $r_1$ or
$r_2r_2r_2$.

One straw-man approach would be: construct DFA out of
the above regular expression, and then the zk-proof for discharging
$s$ would be simply linear with $|s|$, given that
encoding finite state machine is ``free" using pre-processed
lookup arguments such as $\cq$ \cite{cq}. Unfortunately,
it is known that the worst complexity of constructing DFA
out of regex is exponential, and so is negation.
% check paper_data/dfa_size/main_dfa.dat
% gen_regex: 155 out of 3479, pcre: 77 out of 550(to be updated),
% pm_reg: 16067 out of 34279
% total: 16.3k out of 38k 
% ==== SUMMARY DFA Stats =====
%Failed: 16296 Success: 22001
%Subsis (Success): 109274 AVG DFA size: 81.52542233285136
We performed an empirical (best-effort) attempt to convert
all sub-signatures of ClamAV set to DFA, using a server with 128GB
ram (running with 32-threads). A failure is reported when the
conversion times out at 100 seconds, or exhausts the 128GB RAM.
Out of the 38k logical signatures, there are 16.3k failed to
generate DFA. 

As an example, observe sub-signature $1$ in the following
virus signature.
Setting the repetitions 
(replacing the \texttt{\{28\}}) to $2$ results in a minimized
DFA of $1622$ states, but at $16$ it already has $56910$ states.
The NFA to minimized DFA conversion crashes at repetition $17$
by exhausting all $128GB$ RAM resources.
\vspace*{-0.01in}
\begin{small}
\begin{Verbatim}[breaklines=true, breakanywhere=true]
Win.Trojan.Zeroaccess-6932503-0;Engine:81-255,Target:1;0&1&2;4d535642564d{2}2e444c4c::i;56423521f01f{28}0a00{16}00f0300000ffffff08000000010000000000;61d5893c82478645bf6059f9d3408757
\end{Verbatim}
\end{small}

Combining all, In Figure\ \ref{fig:lcreg}
we present the formal syntax of the regex (covering
the special logical structure and counter constraints) of ClamAV.
This also subsumes the syntax for the PCRE regex used in both
ClamAV and Snort. We call it {\em LC-Regex}. We use $\Sigma$ to
denote an alphabet of 16 hex digits in this paper. 


Here a sub-signature is a full-fledged PCRE excepts
look-around and back-references, as noted earlier the rare
instances are re-written to equivalent standard regex or approximated.  
We skip all ``syntax sugars" such as character classes and repetition operators
as they all can be re-written with the basic LC-Regex.
For example, PCRE regex $\text{/[a-d]/}$ is re-written as
$(61|62|63|64)$.

The intersection and negation is allowed only at the logical
layer, where sub-signatures are composed via logical connectors,
and counter constraints are defined also at this layer.
Counter constraints could be re-written, however, we 
specifically represent them so that approximation techniques
introduced later can leverage the structural information directly.
It is apparent that LC-Regex captures the complete regular language.

\begin{figure}
\begin{small}
\begin{bnf}
$s$  ::=  : \textsf{sub-signature}
| $a$: \textsf{character} $a \in \Sigma$
| $.$: \textsf{any}
| $s*$ | $s+$: \textsf{repetition}
| $s_1 s_2$:  \textsf{concatenation}
| $s_1 \vert s_2$:  \textsf{alternation}
| $(s)$:  \textsf{grouping}
;;
;;
$n$ ::= $[0-9]+$ : \textsf{number}
;;
;;
$r$ ::=  : \textsf{LC-Regex}
| $s$: \textsf{sub-signature}
| $r_1 \vert r_2$: \textsf{alternation}
| $r_1 \& r_2$: \textsf{intersection}
| $\overline{r}$: \textsf{negation}
| $s~ (>\vert=\vert<)~ n$: \textsf{counter constraint}
;;
;;
\end{bnf}
\end{small}
\caption{Regular Expression with Logical Structures and Counter 
Constraints (LC-Regex)} 
\label{fig:lcreg}
\end{figure}

\subsection{Approximations}
Our goal is to design an efficient zk-proof scheme that can batch discharge
an input string against a large collection of regex patterns. We would like
to avoid NFA encoding of regex patterns (to avoid the complexity of
$O(|s||r|)$). Due to the difficulty of constructing DFA, we would not
construct DFA whenever possible (in fact, we show that for the entire
ClamAV rule set, it sufficies to generate DFA for only 6 rules). 
In this section, we present four approximation approaches that allow
us to avoid expensive single NFA discharging approach completely. 
In terms of the cost for batch discharging, their cost increases
in ascending order. Thus, the strategy is to filter all patterns
through these approximations one by one. They are: (1) critical pattern
(denoted as \texttt{CP}), (2) bag of critical pattern sets(\texttt{BW}),
(3) single thread distance encoding (\texttt{STD}), and
DFA (\texttt{DFA}). We describe the idea of each approximation
briefly and present data to support our argument.

\subsubsection{Critical Pattern (\texttt{CP})}
Consider discharging $\texttt{51525354}$ for 
a subsignature $s_1 := \texttt{3132\{2-\}(61626364|41424344)}$.
A human reader can immediately
verify the case without even running the DFA of the pattern,
because none of the long pattern words such as $\texttt{61626364}$
or $\texttt{41424344}$ appears in the target string.
These long pattern words are required for any string to
match the subsignature. We call them ``critical patterns".
Formally it is defined as follows:

\begin{definition} \label{deff:cp}
Given a regular expression $r$, a set of strings
$\cp(r)$ is called a {\em critical pattern set} for $r$
if for any string $t\in r$ there exists a string $s \in \cp(r)$
s.t. $s$ is a substring in $t$.
\end{definition}

Intuitively, if none of the words in $\cp(r)$ appears in a string
$s$, then $s$ can be discharged for $r$.  One should note that a 
regular expression may have many
critical pattern sets.  For instance for $s_1$
mentioned above, both $\{\texttt{61626364,41424344}\}$
and $\{3132\}$ are valid \cp, however, $\{\texttt{61626364}\}$
is not, because $\text{313241424344}$ is a valid match for
$s_1$ and it does not have to contain $\texttt{61626364}$.  
In practice, we use a heuristic to collect
those that are the longest and those appear less frequently 
in other subsignatures (thus $\{3132\}$ is not preferred for $s_1$). 
Later, we show that \cp can be constructed from \bw.

Consider a simple scenario that all regular patterns has only
one hex subsignature and no counter constraints.
Given a large collection of regular patterns $R = \{r_1, ..., r_n\}$,
we collect one critical pattern for each $r_i$, and construct
an Acho-Corasick autaton (AC-DFA) on all these fixed word patterns.
Notice that AC-DFA is deterministic, and its
construction complexity is superior compared with constructing DFA
from regex: it is {\em linear} in the total length of fixed words.
Then given an input string $s$, we can simply collect the output
words of AC-DFA along its acceptance path. Those regular patterns
that do not appear (to be associated) with the collected output words
do not have to be considered. 
% support data: paper_data/raw_data/main_exec.dat
% max patterns 198 needed after critical pattern
% (38297-198)/38297
In this case, we can discharge over $99.48\%$ ClamAV rules 
for the vast majority of binary executables and email files 
in Enron data set (in the worst case there are $198$ out of 38297 
ClamAV rules left to discharge for the next approximation, 
for one large executable file. For other executables and
emails in Enron data set, the number of left-over is much smaller, as
their size is smaller and are less likely to contain a lot of
critical patterns).

In practice, we need to address more delicacies. For instance,
we need to construct two ACDFAs, one for case sensitive
and one for case-insensitive. We also need to consider
the logical structure of a ClamAV rule which involves multiple
subsignatures (and their negation cases). For instance,
subsignatures that are negated and counter constraints such as
$0<2$ do not have critical pattern set (that is, they cannot be discharged
via critical patterns). Also, we need to consider the logical structures.
For instance, for a logical signature $(r_1 \vert r_2) \wedge (r_3 \vert r_4)$
we only need to pick one out of $\cp(r_1) \cup \cp(r_2)$ and
$\cp(r_3) \cup \cp(r_4)$.


\subsubsection{Bag of Critical Pattern Set (\bw)}
Bag set can be regarded as a straightforward extension of
the critical pattern set model. It is a {\em conjunction} of
a set of critial pattern sets.  To discharge via
$\bw$ in zk-prover, however, is rather different 
from the previous approach. We rely on ACDFA to count word occurences
and a ternary logic system \cite{ternarylogic},
which allows reasoning about ``maybe".

A ternary-value logic system has three values in the system:
\true ~(True), \false ~(False), and \unknown ~(Unknown). 
We need three operators: $\neg$ (negation), $\wedge$ (and),
and $\vee$ (or), and their truth tables are defined in Figure
\ \ref{fig:ternary}.

\begin{figure}
\begin{small}
\begin{displaymath}
\begin{array}{c|c|c|c}
	\wedge & \true & \false & \unknown \\
	\hline
	\true & \true & \false & \unknown \\
	\unknown & \unknown & \false & \unknown \\
	\false & \false & \false & \false 
\end{array}
\qquad
\begin{array}{c|c|c|c}
	\vee & \true & \false & \unknown \\
	\hline
	\true & \true & \true & \true \\
	\unknown & \true & \unknown & \unknown \\
	\false & \true & \false & \unknown
\end{array}
\qquad
\begin{array}{c|c|c|c}
	\neg & \true & \false & \unknown \\
	\hline
		 & \false& \true & \unknown \\
\end{array}
\end{displaymath}
\end{small}
\caption{Truth Table of Ternary-Logic}
\label{fig:ternary}
\end{figure}


\begin{definition} \label{deff:bw} Given a regular expression
subsignature $r$, a bag of critical pattern set for $r$,denoted as
$\bw(r)$, is a set s.t.  each of its element is a distinct $\cp(r)$.  
\end{definition}


Intuitively, a word $s \in r$ must
satisfy each critical pattern set in $\bw(r)$.
In another word,
if there exists $\vec{u} \in \bw(r)$ s.t. 
for all $t\in \vec{u}$: $t$ is not a substring
of $s$, then $s \not\in r$.

\begin{example} \label{examp:bw} 
Given a regex $r := \texttt{6162(4142|4344)5152}$,
both $\texttt{\{\{6162\}$, $\{4142,4344\}$, $\{5152\}}\}$ and
$\texttt{\{\{6162\}, \{5152\}\}}$ are $\bw(r)$.
\end{example}


In Figure\ \ref{fig:bw} we present algorithms for generating $\bw$,
and for discharging a string by $\bw$. 
$\texttt{gen\_bag\_sets()}$
is defined based on the recursive definition of a subsignature.
One interesting point is that dropping a critical pattern set 
from a $\bw(r)$ yields a valid but coarser $\bw(r)$,
as shown in Example\ \ref{examp:bw}. 
When handling the wildcard \texttt{.} case, the true $\bw$ 
is $\{\Sigma\}$, but we return $\emptyset$ to avoid state explosion.
Another interesting point is how alternation is handled. For instance,
for regex subsignature \texttt{ab((12.*34)|(56.*78))cd} the corresponding 
bag of sets generated are $\{\{\texttt{ab}\},$ $\{\texttt{12}, \texttt{56}\},$ 
$\{\texttt{12}, \texttt{78}\},$ $\{\texttt{34}, \texttt{56}\},$ 
$\{\texttt{34}, \texttt{78}\}, \{\texttt{cd}\}\}$,
resulting from a cross product operation of 
$\bw(\texttt{12.*34})$ and
$\bw(\texttt{56.*78})$. 
In practice, cross-product may result in explosion of terms.
We use heuristics to drop terms above a certain pre-set limit (e.g., $2048$). 
Such an approximation is conservative.

Before we discuss how to discharge in batches using $\bw$, 
consider how ACDFA can be leveraged first. 
Given a large collection of subsignatures $\{r_1, ..., r_n\}$,
an ACDFA can be constructed from $\bigcup_{i\in[n]} \bw(r_i)$.
Then running a string $s$ through the ACDFA, we obtain
an acceptance path, from which, one can obtain a hashmap $\rho$ which
maps from each accepted word to the {\em count} of the word in the string.
Notice that as all patterns are long (the vast majority has length
greater than 10), the ratio of final states 
in the acceptance path is low, and the size of $\rho$
is much smaller than $|s|$ itself. {\bf In our data set (full
ClamAV signatures over 700MB Linux binary executables),
% see paper_data/raw_data/discharge/main_exec.txt (last line)
$|\rho|$ is only $0.06\%$ of $|s|$,
It implies a speed up factor of over $1600\times$ compared with
discharging using DFA directly.}

With such information, we can evaluate counter constraints efficiently,
and evaluate a regex subsignature is a special case of it.
Observe function $\texttt{eval\_cs\_bs}(\rho, r, op, n)$
in Figure\ \ref{fig:bw}.  The function takes the following input:
word frequency map $\rho$, a counter constraint $r ~op~ n$, and 
it tries to decide if the string represented by $\rho$ can be evaluated 
by $\bw$. It returns a ternary logic value, and when
it returns $\true$ ($\false$), it implies that the string generating $\rho$
really meets (does not meet) the counter constraint. A $\unknown$
indicates an unsure value. Thus the approach is conservative.
Then, the ternary value allows composing the evaluation of
subsignatures from an LC-regex, as shown in function
$\texttt{eval\_lc\_bs}()$.

\def\gbw{\texttt{gbs}}
\def\hp{\hspace*{0.1in}}
\begin{small}
\begin{algorithm}[bt]
\SetKwData{Left}{left}\SetKwData{This}{this}\SetKwData{Up}{up}
\SetKwFunction{Union}{Union}\SetKwFunction{FindCompress}{FindCompress}
$\bw \leftarrow {\texttt{\textbf{gen\_bag\_sets}}}(s)$: 
\newline
\# Generate $\bw$ for subsignature $s$ \newline
(1) Denote \texttt{gen\_bag\_sets} as \gbw. \newline
(2) Return \bw(s) based on its structure:
\[
\left\{ 
\begin{array}{l l l}
	a \in \Sigma^+ &\Rightarrow & \{\{a\}\} \\
	. 	& \Rightarrow & \emptyset \texttt{~\# to avoid state explosion} \\
	s_1s_2 & \Rightarrow& \gbw(s_1) \cup \gbw(s_2) \\
	s_1 *& \Rightarrow& \emptyset \\
	s_1+ & \Rightarrow& \gbw(s_1) \\
	s_1~|~s_2 & \Rightarrow& \{\vec{u} \cup \vec{v}~|~ \vec{u} \in \gbw(s_1) \wedge \vec{v}\in \gbw(s_2) \}
\end{array}
\right.
\]

$\true/\unknown/\false \leftarrow \texttt{\textbf{eval\_counter\_cs\_bs}}(\rho, r, \op, n)$:  \newline
\# Evaluate counter constraint by bag of critical pattern sets \newline
(1) Compute $\occ = \min\left(\left\{ \sum_{w \in \vec{v}} \rho(w) ~|~ \vec{v} \in \bw\right\}\right)$. \newline
(2) \return ternary value according to the following chart:
\[
\begin{array}{c|c|c|c}
	\op & \true & \false & \unknown \\
	\hline
      $>$  & - 		& $\occ <=n$ & $\occ > n$ \\
	\hline
      $<$  & $occ<n$ & - 		 & $\occ > n-1$ \\
	\hline
      $= (\texttt{where} n>0) $  & - & $\occ <n$ & $\occ > n-1$ \\
	\hline
      $= (\texttt{where} n=0) $  & \occ=0 &  -  & $\occ > 0$ \\
\end{array}
\]

$\true/\unknown/\false \leftarrow \texttt{\textbf{eval\_subsig\_bs}}(\rho, r)$:  \newline
\# Evaluate subsignature by $\bw$ \newline
\return $\texttt{eval\_counter\_cs\_bs}(\rho, r, >, 0)$

\def\dlb{\texttt{dlb}}
$\true/\unknown/\false \leftarrow \texttt{\textbf{eval\_lc\_bs}}(\rho, r)$:  \newline
\# Evaluate LG-regex by $\bw$ \newline
\# Here $r_i$: LC-regex, $s_i$: subsignature \newline
(1) Denote \texttt{eval\_lc\_bs} by texttt{dlb}.\newline
(2) \return one of the following based on the structure of $r$:
\begin{footnotesize}
\[
\left\{
\begin{array}{l l l}
	s & \Rightarrow & \texttt{eval\_subsig\_bs}(\rho, s) \\
	s ~\op~ n & \Rightarrow & \texttt{eval\_counter\_cs\_bs}(\rho, s, \op, n) \\
	\overline{r_1} &\Rightarrow & \neg \texttt{eval\_lc\_bs}(\rho,r_1) \\
	r_1 | r_2  &\Rightarrow &
			\texttt{eval\_lc\_bs}(\rho,r_1) \vee 
			\texttt{eval\_lc\_bs}(\rho,r_2)  \\
	r_1 \& r_2 & \Rightarrow &
			\texttt{eval\_lc\_bs}(\rho,r_1) \wedge 
			\texttt{eval\_lc\_bs}(\rho,r_2) 
\end{array}
\right.
\]
\end{footnotesize}
\caption{Bag of Critical Pattern Set Algorithms}
\label{fig:bw}
\end{algorithm}
\end{small}

We now delve into the details of $\texttt{eval\_counter\_cs\_bs}()$
by an example. Given counter constraint $r ~\op~ n$, it first
computes $\bw(r)$. For each critical pattern set in $\bw(r)$,
it uses $\rho$ to compute the sum of occurences of all the words
in the critical set pattern, and denote that as the occurrence for
that pattern set. Then it iterates through all the pattern sets
and get the minimum occurrence (denoted as $\occ$ in the function).
It then use $\occ$ to decide if the string behind $\rho$ satisfies
the constraint.

\begin{example} 
Let $r := \texttt{aa(11|22)bb}$, and let the
counter constriant be $r < 2$. 
One valid $\bw(r)$ is 
$\{\{\texttt{aa}\},$ $\{\texttt{11,22}\},$ $\{\texttt{bb}\}\}$.
Given string $s := \texttt{aa112211bbaa}$, its word frequency
map $\rho$ is: 
$\{\texttt{aa} \rightarrow 2$,
$\texttt{11} \rightarrow 2$,
$\texttt{22} \rightarrow 1$,
$\texttt{bb} \rightarrow 1\}$. Given $\rho$, the algorithm
can compute the total of occurrence for each of the three
critical pattern sets. They are $2$, $3$, and $1$, respectively
for $\{\texttt{aa}\},$ $\{\texttt{11,22}\},$ $\{\texttt{bb}\}$.
Then $\occ$ is set to minimum of them, i.e., 1.
Following the chart, it decides $s$ satisfies the counter constraint
$r < 2$. The intuition is clear: the required pattern
$\texttt{bb}$ only appeared once, there is no way for the entire string
to have more than $2$ matches of $r$.

Fix $r := \texttt{aa(11|22)bb}$. The following tables lists more examples.
A string $s$ is discharged of a regular subsignature $r$ when
the output is $\false$ (a sure false value in ternary logic).

\begin{small}
\begin{center}
\begin{tabular}{c | c | c | c}
\hline
$s$ & \texttt{counter constraint} & $\occ$ & output \\
\hline
\hline
\texttt{1122bbaabb} & $r>0$ & 1 & $\unknown$ \\
\texttt{11aa} & $r>0$ & 0 & $\false$ \\
\texttt{aa11bbaa22bb} & $r=2$ & 2 & $\unknown$ \\
\texttt{aa22bbbbbbbb} & $r<2$ & 1 & $\true$ \\
\hline
\end{tabular}
\end{center}
\end{small}

\end{example}

\subsubsection{Single-Thread Distance Encoding ($\sde$)}
We now introduce Single-Thread Distance Encoding ($\sde$) for 
capturing the distance between pattern words, which $\bw$ cannot. 

Consider $r := \texttt{ab(12|34)cd}$, one valid $\sde(r)$ is
$\texttt{ab.\{2,2\}cd}$. It captures the fact that the distance between
$\texttt{ab}$ and $\texttt{cd}$ should be $2$. The motivation
for abstracting away alternation and nested structures in regex 
is based on producing an efficient zk-prover algorithm using ACDFA.
We provide a formal definition of $\sde$ below.

\begin{definition} \label{deff:std} Given a regex subsignature $s$,
we call another regex $r$ a valid Single Thread Distance Encoding
of $s$, written as $\sde(s)$, if $L(s) \subseteq L(r)$ and $r$
satisfies the following BNF syntax: 
\[
r := a \in \Sigma^+ ~|~ .\{n_1, n_2\} ~|~ r_1 r_2,
\]
where $r_1$ and $r_2$ are $\sde$ of some regex subsignatures.
Here $n_1 \leq n_2$ and we abuse the notation of $n_2$
to allow $\infty$, so $.+$ can be represented as
$.\{1,\infty\}$.
\end{definition}

Each $\sde$ regex has a unique simplified
$\sde$ which consists of alternating wildcard strings and
fixed patterns. For instance $\texttt{.\{2,5\}ab.\{3,3\}.\{4,5\}cd}$
can be simplified to $\texttt{.\{2,5\}ab.\{7,8\}cd}$.
The function $\texttt{simplify}(s)$ is given in Fig\ \ref{fig:sde}.
It applies a reduction of merging consecutive fixed words and
wildcards and produces an equivalent $\sde$. Note that 
$\concat$ stands for fixed word concatenation in its function body.

Similarly to $\bw$, we can define $\sde$ recursively
based on the subsignature structure. 
This is implemented in function $\texttt{gen\_subsig\_sde}(s)$ in
Figure\ \ref{fig:sde}. It relies on the $\texttt{wildcard}(s)$
function which approximates a 
a fixed word to wildcards, e.g.,
from $\texttt{abc}$ to $\texttt{.\{3,3\}}$.
Then using it, alternations are abstracted away
by keeping track of the length bounds of all components,
e.g.,
$\texttt{(ab|cdef)}$ is approximated to
$\texttt{.\{2,5\}}$. It is clear that the
approximation is conservative in the sense that
any fixed word $w \in L(s)$ also belongs to $L(\sde(s))$.

In function $\texttt{eval\_counter\_cs\_sde}(w, s, \op, n)$
we evaluate whether fixed word $w$ satisfies counter
constraint $s~\op~n$, and evaluation of a general subsignature
is a special case of it. Then using the ternary logic
one can extend to LC-regex.

We present an example to illustrate how the algorithm works.
Again, like the $\bw$ method, the approach relies on
ACDFA to generate a map $\rho$ which maps each accepted
word to a set of indices where it appears in the target string.
When the function is called, $\rho$ instead of
the target string is passed for mass processing. 

In practice, $\rho$ needs to go through two levels
of table projections to optimize zk-prover performance.
For the first projection, after the $\bw$ completes,
the zk-prover collects all signatures that need to go to
$\sde$, and project $\rho$ to contain entry words that
are related to signatures failed by \bw only.  
This is a one time projection and the cost is
$O(|s|)$. 
Let there be $u$ signatures to handle in the
$\sde$ stage, usually $u>100$ for file size $2^{20}$.
Given the cost is spread to all signatures, the first
projection cost is negligible in average for discharging a signature
via $\sde$.
The first projection results in a hashmap
$\rho^{(1)}$,
% check end of paper_data/raw_data/discharge_data/main_exec.dat
%entry: 
which is about $5\%$ of $|s|$.\footnote{The result is obtained 
at bounding the $\sde$ fixed patterns with length at least $4$.
It is possible to achieve better compression ratio by increasing
minimum length, however, reducing effectiveness of $\sde$ thus
leaving more signatures to the DFA stage.}

To further cut the cost, for each signature involved (let it be written $r_i$),
$\rho^{(1)}$ is further projected into a $\rho^{(2)}_i$,
an individual table for $r_i$.  
Layer 2 projection can be processed in batch, as a
variation of permutation argument for an inflated table from $\rho^{(1)}$.
The maximum inflation rate according to our data,
% check end of paper_data/raw_data/discharge_data/main_exec.dat
% entry: 
is only $2.17$. This implies the total combined cost for
layer 2 table projection is  $O(|\rho^{(1)}|)$ in total for all signatures. 

Then $\rho^{(2)}_i$ is used for the
$\sde$ process for each signature $r_i$, note that
$\sum_{i=1}^u |\rho^{(2)}| \approx |\rho^{(1)}|$. 
The combined $\sde$ cost for all subsignatures is thus $O(\rho^{(1)}$.
The concrete data shows that the combined $\sde$ witness length
% pm_reg (sde) total witness_len/file_size: (avg: max)
in average is only $2.5\%$ (and max $60.8\%$)
of the file size.
Thus, in total for $u$ signatures, the zk-prover pays
combined cost only a constant factor (of $<2$) of file size.
Compared with DFA discharging, the speed-up of $\sde$
depends on the value of $u$, which is typically $>100$ for file size $2^{20}$.

\begin{example} Let $r = \texttt{ab(12|32)cd}$ and
$w = \texttt{ab11cd00ab22cd}$. The goal is to
evaluate $w$ against counter constraint $r>2$. 

The algorithm first collects bag of words from all related
signatures (which includes $\texttt{ab}$ and $\texttt{cd}$).
Then it computes 
$\rho$ for $w$ with ACDFA, where $\rho = \{
	\texttt{ab} \rightarrow \{0, 8\},
	\texttt{cd} \rightarrow \{4, 12\},
\}$, e.g., $\texttt{ab}$ appeared at indices $0$ and $8$
in $w$.

When processing
$\texttt{eval\_counter\_cs\_sde}(\rho, r, >, 2)$,
the algorithm first computes repeated $r$ as the following,
to require $3$ non-overlapping match of $r$:
\[
\texttt{.*ab(12|32)cd.*ab(12|32)cd.*ab(12|32)cd}.*
\]
Its simplified $\sde$ is:

\[
\{0,\infty\}
\texttt{ab}.\{2,2\}
\texttt{cd}.\{0,\infty\}
\texttt{ab}.\{2,2\}
\texttt{cd}.\{0,\infty\}
\texttt{ab}.\{2,2\}
\texttt{cd}.\{0,\infty\}
\]
Note that we have $6$ fixed pattern words in the above $\sde$.
Using $\rho$, the allowed position set are computed, as shown
below:
\begin{center}
\begin{tabular}{c | c | c | c }
\hline
 $\varphi_0 = \{0\}$  &
 $\varphi_1 = \{0, 8\}$  &
 $\varphi_2 = \{4, 12\}$  &
 $\varphi_3 = \{8\}$  
 \\
\hline
 $\varphi_4 = \{12\}$  &
 $\varphi_5 = \emptyset$  &
 $\varphi_6 = \emptyset$  &
 - \\
\hline
\end{tabular}
\end{center}

Intuitively, $\varphi_1$ reflects the permitted positions
of the 1st occurrence of \texttt{ab}, $\varphi_2$ the positions
for the 1st occurrence of \texttt{cd}, etc. Each round, the algorithm
proceeds from the the last position set, advances with the length
of the last word and the wildcard length bound in between, and
then intersect with with the set of occurrence indices of the
next word. 

As shown in the computation, there is no valid position
for the third pair of $\text{ab}$ and $\text{cd}$.
Following the chart, ternary value $\false$ is returned,
which discharges $w$ for counter constraint $r>2$.
\end{example}


\begin{small}
\begin{algorithm}[bt]
\SetKwData{Left}{left}\SetKwData{This}{this}\SetKwData{Up}{up}
\SetKwFunction{Union}{Union}\SetKwFunction{FindCompress}{FindCompress}
$\sde \leftarrow {\texttt{\textbf{simplify}}}(s)$: 
\newline
\# produce a simplified version of $s$ where $s$ is a $\sde$ for some $r$ 
\newline
Apply reduction below to $s$ until no more can be applied. \newline
\[
 \left\{
 	\begin{array}{l l l}
		a_1 a_2$ \textrm{ where } $a_i \in \Sigma^+ & \Rightarrow & a_1 \concat a_2 \\
		.\{m_1, m_2\} .\{n_1, n_2\} & \Rightarrow & .\{m_1+n1, m_2+n2\}

	\end{array}
 \right.
\]

$\sde \leftarrow {\texttt{\textbf{wildcard}}}(s)$: 
\newline
\# Assuming $s$ is already a $\sde$, produce its wildcard approximation 
\newline
(1) Let $s'$ be the result of approximating each fixed word $a$ 
in $s$ to $.\{|a|,|a|\}$. \newline
(2) \return $\texttt{simplify}(s')$.

\def\gsde{\texttt{gss}}
$\sde \leftarrow {\texttt{\textbf{gen\_subsig\_sde}}}(s)$: 
\newline
\# Generate $\sde$ for subsignature $s$ \newline
(1) Denote \texttt{gen\_subsig\_sde} as \gsde. \newline
(2) Compute $s_2$ based on the structure of $s$:
\[
\left\{ 
\begin{array}{l l l}
	a \in \Sigma^+ &\Rightarrow & a \\
	. 	& \Rightarrow & . \\
	s_1s_2 & \Rightarrow& \gsde(s_1) \concat \gsde(s_2) \\
	s_1 *& \Rightarrow& \texttt{.\{0,}\infty\texttt{\}} \\
	s_1+ & \Rightarrow& \gsde(s_1) || \texttt{.\{}0,\infty\texttt{\}} \\
	s_1~|~s_2 & \Rightarrow& \texttt{.\{}a,b\texttt{\}} \textrm{ where } \\
	& & \texttt{.\{}u,v\texttt{\}} = \texttt{wildcard}(\gsde(s_1)) 
		\textrm{ and } \\
	& & \texttt{.\{}x,y\texttt{\}} = \texttt{wildcard}(\gsde(s_2))  \\
	& & a = \min(u,x) \textrm{ and } b = \max(v,y)
\end{array}
\right.
\]

$\true/\unknown/\false \leftarrow \texttt{\textbf{eval\_counter\_cs\_std}}(\rho, r, \op, n)$:  \newline
\# Evaluate counter constraint by SDE \newline
\# $\rho[u]$ is the set of $u$'s occurrence indices in target string to eval.
\newline
(1) Let $n' = \left\{
	\begin{array}{l l l}
		r > n & \Rightarrow & n+1 \\
		r = 0 & \Rightarrow & 1 \\
		r = n (n>0) & \Rightarrow & n \\
		r < n & \Rightarrow & n
	\end{array}
	\right.
$. \newline
(2) Define $r'$ as $r$ repeated $n'$ times intermixed
with $\texttt{.*}$. \newline
(3) Compute $e' = \texttt{simplify}(\sde(r'))$.
Write $e'$ as the following:
\[
\texttt{.\{}a_1, b_1\texttt{\}} 
w_1 
\texttt{.\{}a_2, b_2\texttt{\}} 
w_2 
\ldots 
\texttt{.\{}a_k, b_k\texttt{\}} 
w_k
\texttt{.\{}a_{k+1}, b_{k+1}\texttt{\}} 
\]
(4) Define $\varphi_0 = \{0\}$ and $w_0 = \epsilon$.
For $i \in [k]$, iteratively compute:
\[
\varphi_i = \{[u + a_i + |w_{i-1}|, u + b_i + |w_{i-1}]
	\intersect \rho(w_i)
	~|~ u\in \varphi_{i-1}\}
\]
(5) \return ternary value according to the following chart:
\[
\begin{array}{c|c|c|c}
	\op & \true & \false & \unknown \\
	\hline
      >  & - & \varphi_k = \emptyset & \varphi_k \neq \emptyset \\
	\hline
      <  & \varphi_k = \emptyset & - & \varphi_k \neq \emptyset \\
	\hline
      = (\textrm{where}~n>0)  & - & \varphi_k = \emptyset & \varphi_k \neq \emptyset \\
	\hline
      = (\textrm{where}~n=0)  &  \varphi_k = \emptyset & - & \varphi_k \neq \emptyset \\
	\hline
\end{array}
\]


\caption{Single Thread Distance Encoding Algorithms}
\label{fig:sde}
\end{algorithm}
\end{small}

\subsubsection{DFA and Individual $\sde$}
When a signature cannot be discharged by $\sde$, it goes to the $\dfa$ stage
where each subsignature is translated to a deterministic finite state machine,
then the target string is fed to each DFA and then the results are combined
using the logical expression as defined in the signature. Note that 
this is not an approximation, as long as the DFAs can be constructed
for each subsignature within the constraint of time and RAM resources available,
which is unfortunately not the case.

In theory, we can discharge a string using NFA (by enumerating all possible
state set reached by the input string), as done in Zombie \cite{Zombie23}).
However, we would like to avoid the cost of $O(|\texttt{NFA}| \times |s|)$.
In this case, we will run another round of $\sde$.

Recall that in $\sde$ there is a minimum requirement of pattern word length (e.g., $4$
nibbles). The reason is that if the pattern word is too short, it blows up the
size of the final states along the ACDFA acceptance path. There are very few
ClamAV signatures which have short pattern words, meanwhile, its DFA is
too costly to construct, e.g., as shown below, where
the pattern words are single characters intermixed with variable length
wildcards, which fails all aforementioned discharging methods.
In addition, its use of named back-references exceeds capability
of regex.

\begin{small}
\begin{Verbatim}[breaklines=true, breakanywhere=true]
Js.File.MaliciousHeuristic-6260249-0;Engine:81-255,Target:7;0>1&1>7&2&3;6576616c28;76617220;0&1/w(?P<junk>[a-z\d]{3,6})s(?P=junk)c(?P=junk)r(?P=junk)i(?P=junk)p(?P=junk)t(?P=junk).(?P=junk)s(?P=junk)h(?P=junk)e(?P=junk)l(?P=junk)l/; ... 
\end{Verbatim}
\end{small}

In this case, we resolve to {\em Individual ACDFA \sde} (shortened $\isde$).
We extract the bag of words from the signature again, without the minimum word length.
Construct an ACDFA for that subsignature only, and run the $\sde$ approach again.
In this case, the ACDFA does not ``mass process" a lot of signatures, but just one.
This approach tends to be more expensive than DFA, and we place it in the last round.

\subsection{Preliminary Evaluation}
Before we encode all methods in zkSNARK, we evaluate their effectiveness
using real data set. We use the entire ClamAV logical signature set,
which comprises of $38876$ signatures.\footnote{Thirteen signatures are dropped
from the original $38889$ signatures because they report false positives on
some Linux executables or Enron emails which are benign. 
By ``benign" we mean: when running clamav using the specific rule set,
ClamAV does not report positive on matching virus signatures.
We do not handle the \texttt{Target} meta-data, i.e.,
signatures are applied to all files. In this case, 
\texttt{Win.Trojan.FakeSkype}, which was targeted at Win32 PE,
raises false positive \texttt{vmlinuz} (Linux kernel image). 
The rest have dependency on program entry 
point position,
e.g., 
\texttt{Win.Trojan.Generic-43;...;EP+0:5250{-8}5133c0;...}.
As the position constraint needs to locate executable file entry point,
which exceeds regex ability, we replace \texttt{EP+0:} with wildcard string
\texttt{.*} before the patterns. However, such patterns are too short and
cause false positives. In this case, 8 signatures are removed from the set.
We use multiline mode for all signatures, i.e., new line character is
matched by wildcard dot. In this case, one signature (which works in
single line mode) is removed
for causing false positive due to multi-line approximation.}
Among them, the vast majority have hex subsignatures and $559$ of them have
PCRE subsignatures. Very few ($6$) of these PCRE signatures use features such
as named back-references and look-arounds, and they are approximated.

We use two sets of target data set (to discharge free of malware). The first 
is a set of all binary executable files in a Linux CentOS ($2702$ files
total $752$ MB) and Enron Email Data Set ($517401$ files totally $2.74$ GB).
Running a single thread clamav to discharge all binary executables takes about
$10$ seconds, but much longer ($4$ hours) on Enron Data 
set (mainly due to file reading cost).

\begin{figure}
\begin{footnotesize}
\begin{center}
\begin{tabular}{c c c c c c c}
\hline
\multicolumn{7}{c}{\bf ClamAV Logical Signatures (38886)} \\
\hline
\multicolumn{7}{c}{Centos Binary Executables (2702 samples)} \\
\hline
log(file) & files & After $\cp \Rightarrow$ & $\bw \Rightarrow$ & $\sde \Rightarrow$ & 
	$\dfa \Rightarrow$ & $\isde \Rightarrow$ \\
\hline
\hline
7 &   72     & (87,90,72)  & (87,87,72)  & (0,0,0) & (0,0,0) & (0,0,0)\\
8 &   0  & (0,0,0) & (0,0,0) & (0,0,0) & (0,0,0) & (0,0,0)\\
9 &   0  & (0,0,0) & (0,0,0) & (0,0,0) & (0,0,0) & (0,0,0)\\
10&   2  & (96,101,2)  & (90,93,2)   & (0,0,0) & (0,0,0) & (0,0,0)\\
11&   13     & (96,103,13) & (89,93,13)  & (0,0,0) & (0,0,0) & (0,0,0)\\
12&   16     & (98,105,16) & (90,93,16)  & (0,0,0) & (0,0,0) & (0,0,0)\\
13&   105    & (98,105,105)    & (90,94,105) & (0,0,0) & (0,0,0) & (0,0,0)\\
14&   727    & (105,119,727)   & (92,97,727) & (0,1,2) & (0,0,0) & (0,0,0)\\
15&   478    & (104,162,478)   & (92,138,478)    & (0,2,1) & (0,0,0) & (0,0,0)\\
16&   442    & (107,163,442)   & (92,139,442)    & (0,2,70)    & (0,0,0) & (0,0,0)\\
17&   273    & (116,164,273)   & (94,137,273)    & (0,2,6) & (0,0,0) & (0,0,0)\\
18&   163    & (122,183,163)   & (95,144,163)    & (0,2,7) & (0,0,0) & (0,0,0)\\
19&   158    & (131,191,158)   & (96,142,158)    & (0,2,60)    & (0,0,0) & (0,0,0)\\
20&   67     & (145,201,67)    & (99,147,67) & (1,3,55)    & (0,0,0) & (0,0,0)\\
21&   153    & (185,202,153)   & (132,149,153)   & (1,3,146)   & (0,0,0) & (0,0,0)\\
22&   14     & (150,173,14)    & (99,103,14) & (1,3,11)    & (0,0,0) & (0,0,0)\\
23&   6  & (175,235,6) & (114,148,6) & (3,4,6) & (0,0,0) & (0,0,0)\\
24&   9  & (224,277,9) & (139,174,9) & (2,4,9) & (0,1,1) & (0,0,0)\\
25 &  3  & (226,270,3) & (146,168,3) & (3,4,3) & (0,1,2) & (0,0,0)\\
26 &  0  & (0,0,0) & (0,0,0) & (0,0,0) & (0,0,0) & (0,0,0)\\
27 &  1  & (285,285,1) & (184,184,1) & (4,4,1) & (0,0,0) & (0,0,0)\\
\multicolumn{7}{c}{Enron Email Data (517401 samples)} \\
\hline
9 &   14025  & (87,147,14025)  & (83,139,14025)  & (0,0,0) & (0,0,0) & (0,0,0)\\
10 &   141762     & (88,148,141762) & (84,140,141762) & (0,0,0) & (0,0,0) & (0,0,0)\\
11 &  175664     & (90,156,175664) & (85,144,175664) & (0,0,0) & (0,0,0) & (0,0,0)\\
12 &  113795     & (91,160,113795) & (86,146,113795) & (0,0,0) & (0,0,0) & (0,0,0)\\
13 &  47882  & (93,166,47882)  & (87,147,47882)  & (0,0,0) & (0,0,0) & (0,0,0)\\
14 &  17421  & (96,172,17421)  & (87,148,17421)  & (0,0,0) & (0,0,0) & (0,0,0)\\
15 &  4681   & (98,162,4681)   & (88,148,4681)   & (0,0,0) & (0,0,0) & (0,0,0)\\
16 &  1331   & (101,154,1331)  & (89,136,1331)   & (0,0,0) & (0,0,0) & (0,0,0)\\
17 &  644    & (104,164,644)   & (88,139,644)    & (0,0,0) & (0,0,0) & (0,0,0)\\
18 &  175    & (110,177,175)   & (90,142,175)    & (0,0,0) & (0,0,0) & (0,0,0)\\
19 &  12     & (113,119,12)    & (89,91,12)  & (0,0,0) & (0,0,0) & (0,0,0)\\
20 &  3  & (107,108,3) & (91,92,3)   & (0,0,0) & (0,0,0) & (0,0,0)\\
21 &  6  & (105,110,6) & (91,92,6)   & (0,0,0) & (0,0,0) & (0,0,0)\\
\end{tabular}
\end{center}
\end{footnotesize}
\caption{Preliminary Data}
\footnotesize{Column 1: the log of file size. Column 2: the number of file samples.
For the rest of columns, each tuple (avg, max, count) represents the average and max
number of signatures that can not be discharged by the approach. The count represents
the number of files that need to go to next discharging stage.}
\label{fig:predata}
\end{figure}

Given the cost of the proposed schemes, we discharge an input file following
the order of $\cp$, $\bw$, $\sde$, $\dfa$, and $\isde$. Once a signature
is discharged, it does not go to the next stage. At the last stage
a count $0$ means that a file is completely discharged of all virus signatures
(thus malware free).  We summarize the asymptotic complexity of
these methods.  Let $u$ be the number of signatures to discharge
at a stage, their prover complexity are:
$O(|s|)$, $O(|s|)$,  $O(|s|)$, $O(u \times |s|)$, $O(u \times |s|)$, respectively.
Note that for the first $3$ stages, the prover complexity is
{\em independent} of the number of signatures to discharge.
In our real data set, the $u$ for the last two stages is less than $2$.


In Figure\ \ref{fig:predata} we display the effectiveness of the proposed
$5$-stage discharging scheme, over the binary executable data set and
Enron data set. The Enron set performs similarly, while it relies more
on $\sde$ because Enron dataset are mostly alpha-numerical and there are 
many more ``false positives" on patterns if no checks on structures are
performed.

Take the $2^{20}$ file size category in CentOS dataset as an example:
\[
\begin{footnotesize}
\begin{tabular}{c c c c c c c}
log(file) & files & After $\cp \Rightarrow$ & $\bw \Rightarrow$ & $\sde \Rightarrow$ & 
	$\dfa \Rightarrow$ & $\isde \Rightarrow$ \\
\hline
20  &   67       &  (143,198,67)      &  (104,150,67)      &  (1,2,54)      &  (0,0,0)   &  (0,0,0)\\
\end{tabular}
\end{footnotesize}
\]

There are $67$ files in the category of $1MB$ size. The first $\cp$ stages
screens the vast majority ($99.5\%$) of signatures and in average there are
$143$ (max $198$) out of $38k$ signatures left. The next $\bw$ stages
discharged in average $39$ signatures, and the $\sde$ discharged all of
the rest except $1$ (in average) and $2$ (max) for $54$ files. 
None of the signatures for $1MB$ files would have to go through to the
last $\isde$ stage. For the entire CentOS dataset,
in the worst case, we only have
to discharge $7$ signatures via $\dfa$ for one file, however, this happens
very rarely. For both the CentOS and Enron dataset,
there are only two signatures\footnote{$\texttt{Js.File.MaliciousHeurstic-6260249-0}$ and $\texttt{Html.Exploit.CVE\_2017\_0069}$ $\texttt{-6013012-3}$}
that have to to through $\isde$ for $5$ (out of $2702$) binary executable files
and $6$ (out of $517401$) Enron email files.
In summary, given that for the last two stages the number of signatures is less than
$2$, the total prover cost is just a constant factor ($\leq 5$) of the
input string size, i.e.,
\[
	5 \times |s|.
\]
This is a speed-up of $7.7\times 10^{6}$
compared with the state of the art (\cite{Zombie23,Luo23})\footnote{In the case of Reef \cite{Reef} it works very well for the non-membership proof case 
where the patterns start
at fixed positions. However, if there is a \texttt{.*} preceding regex patterns,
our technique works better. It is not known if Reef's alternating automata 
can be extended easily to mass process large collection of signatures.}, which
use NFA discharging each
signature separately, where the average $|\texttt{NFA}|>1000$
(given average length of ClamAV signature is $300$):
\[
38878 \times |s| \times |\texttt{NFA}|.
\]

The speed-up comes from two sources. First, we ``mass-process" large collections
of signatures using ACDFA in the first three stages. Second, we rely on
pre-processed lookup arguments (whose prover complexity is independent
of the lookup/container table size) in proving DFA acceptance, thus eliminating
the size of automata in the prover cost. Several different types of
look-up arguments in used in our zk-prover scheme, which is demonstrated
in details in Section 4.


% \subsection{PM-Reg (TO REMOVE)}
% One interesting observation of the above signature is that
% ClamAV signature set is oriented by long fixed 
% patterns, and this applies to most other malware scanner and
% intrusion detection signature set. From a logical signature,
% one can derive a ``critical fixed pattern" set of the signature,
% so that if all of them do not appear in a word $s$, then
% it is guaranteed that $s$ is not an instance of such a signature.
% For instance, one ``critical fixed pattern" set for
% the \ttt{Fallchill} signature would be:
% \begin{small}
% \begin{Verbatim}[breaklines=true, breakanywhere=true]
% {B8DBFFFFFF412A0341880349FFC341803B0075} 
% \end{Verbatim}
% \end{small}
% That is: missing all elements of the critical fixed pattern set
% meaning sub-signatures could be matched. One could see that
% there can be multiple criticial fixed pattern set for the same signature,
% Here one may take the longest fixed pattern from each related sub-signature.
% We then can design a 2-stage proof system to filter out
% those non-related signatures.
% 
% 
% An initial data anlysis shows that for the full 38k ClamAV signatures,
% and for all 2702 binary executables in CentOS 7, the average
% number of near miss of an executable is only $4$ per executable (that is
% only $4$ out of $38k$ signatures should be checked). The number of 
% near misses grows apparently with file size. It remains an interesting
% problem on if our technique can be developped for chunking a big file
% into smaller sizes and results be aggregated. The following table
% displays the break-down of near miss information.
% For all other $363$ regex patterns, the near-miss approach work,
% but not as well as the PM-reg patterns. The average near
% %see data/paper_data/crit_general.dat
% miss for all 2704 executables is $14$, and the max is $48$.
% 
% \begin{table}
% \begin{tabular}{|c|c|c|c|}
% \hline
% log2(SIZE) &    Count &  Avg & Max \\
% \hline
% 1 &  0 & - & - \\
% 2 &  0 & - &  - \\
% 3 &  0 & - &- \\
% 4 &  0 & - &- \\
% 5 &  0 & - &- \\
% 6 &  0 & - &- \\
% 7 &  72&  0 &  0 \\
% 8 &  0 &  - & - \\
% 9 &  0 &  - & - \\
% 10&   2 &  0 &  1 \\
% 11 & 13 & 0  & 2 \\
% 12 & 16 & 0  & 2 \\
% 13 & 105 & 0 &  2 \\
% 14 & 727 & 1 &  9 \\
% 15 & 478 & 1 &  7 \\
% 16 & 442 & 3 &  10 \\
% 17 & 273 & 5 &  15 \\
% 18 & 163 & 7 &  24 \\
% 19 & 158 & 11&  20 \\
% 20 & 67  & 16&  25 \\
% 21 & 153 & 21&  29 \\
% 22 & 14  & 19&  30 \\
% 23 & 6   & 20&  32 \\
% 24 & 9   & 40&  54 \\
% 25 & 3   & 37 & 54 \\
% 26 &  0   & - & - \\
% 27 & 1   & 49 & 49 \\
% \hline
% \end{tabular}
% \caption{Total: 2702 files, Avg: 4 Near Miss}
% \end{table}
% 
% Another interesting observation of the ClamAV signature set is that
% not all regex features are frequently used. Logical operators such as
% \ttt{and}, \ttt{or}, and \ttt{not} (of entire regex) are frequently 
% used, but repetition and Kleen stars are mostly applied
% to wilecard match ``\ttt{.}" only. This naturally leads to a
% fragment called BM-Refex as defined in the following.
% 
% \begin{bnf}
% $t$  ::=  : \textsf{basic term}
% | $s$: $s \in \Sigma$
% | $.*$: \textsf{wildcard Kleen star }
% | $.\{a,b\}$: \textsf{repetition from a to b times}
% ;;
% $s$  ::=  : \textsf{sub-signature}
% | $t$: \textsf{basic term}
% | $s t$: \textsf{concatenation}
% ;;
% $r$ ::=  : \textsf{Pattern Matching Regex (PM-Regex)}
% | $s$: \textsf{subsignature}
% | $r \& s$: \textsf{intersection}
% | $r \mid s$: \textsf{alternation}
% | $\overline{r}$: \textsf{complement}
% \end{bnf}
% 
% Our initial analysis shows that out of $38330$ ClamAV logical
% signatures, $37.9k$ ($98.5\%$) are BM-Regex.
% \footnote{$416$ signatures
% are excluded from the 
% full set are: 559 PCRE regex patterns for engineering convenience
% (will be fixed later), $216$ signatures relying on content
% position (e.g., the 3rd section offset 120), $18$ signatures
% for complex modes (full-word mode), $108$ for basic terms
% containing class characters and alternation, $72$ for
% pattern counters (e.g., certain pattern appearning than $100$ times
% which is costly for us to model). Given average file size $232kb$,
% the average packed acceptance path is $440k$ (i.e., projecting
% acceptance path to final states only), due to some very 
% small patterns. After removing $173$ BM-Regex signatures containing
% small patterns, we end up with $37741$ signatures, but
% resulting in only $6.6k$ packed acceptance path length, which
% is crucial for the performance of arithmetic circuit.
% % see preprocess.dat 
% } 
% 
% The benefit of BM-Regex is that we can again leverage
% the AC-DFA. We collect all fixed word patterns from all BM-Regex
% patterns and this results in a AC-DFA of $6$ million states
% and $96$ million transitions. The size of AC-DFA does not matter
% as the prover cost will be independent of its size once pre-processing
% is done. From an acceptance path of a string, we collect
% a {\em packed}  path consisting accepting states and their positions
% only. Usually this packed path is much smaller than the
% original string size. Our data analysis shows that for all
% CentOS 7 executable files, the pack ratio is $437$. For instance
% for all files of about $1Mb$, the packed trace is average is
% only $3k$ elements. 
% 
% We then have an arithmetic circuit for each signature, which
% encodes the arithmetic constraints that are caused by
% the wildcard match patterns. This can be done very efficiently,
% and checking only on the packed path. For instance, given
% a $1MB$ file, in average, there  $16$ signatures involved,
% in average, the signature pattern involves no more than $5$ components
% (5 sub-signatures). Each arithmetic circuit costs
% $25 \times 3000$ constraints, and $16$ of them only
% costs $700k$ R1CS constraints. The concrete cost is only
% a constant factor of generating the commitment itself.

